% LEGACY DOCUMENT (historical data paper wave 19)
% Current status: See RESEARCH_STATUS.md and TECH_TREE.md for canonical assessment.
% δ_CP prediction is WITHDRAWN (>5σ excluded by NuFIT-6.0 + T2K+NOvA 2025).
% See PREDICTIONS_PREREGISTERED.md for canonical up-to-date assessment.

%%%%%%%%%%%%%%%%%%%%%%%%%%%%%%%%%%%%%%%%%%%%%%%%%%%%%%%%%%%%%%%%%%%%%%%%%%%%%%%
%  H₄ Numerological Atlas: A Catalog of Standard Model Parameter Formulas
%
%  Data paper — PLOS ONE submission
%  Document class: standard article (PLOS ONE accepts any LaTeX)
%  Version: Wave 19
%%%%%%%%%%%%%%%%%%%%%%%%%%%%%%%%%%%%%%%%%%%%%%%%%%%%%%%%%%%%%%%%%%%%%%%%%%%%%%%

\documentclass[11pt,a4paper]{article}

% ─── Packages ─────────────────────────────────────────────────────────────────
\usepackage[utf8]{inputenc}
\usepackage[T1]{fontenc}
\usepackage{amsmath,amssymb,amsthm}
\usepackage{booktabs}
\usepackage{longtable}
\usepackage{graphicx}
\usepackage{hyperref}
\usepackage{url}
\usepackage[numbers,sort&compress]{natbib}
\usepackage{caption}
\usepackage{xcolor}

% ─── Hyperref setup ───────────────────────────────────────────────────────────
\hypersetup{
    colorlinks=true,
    linkcolor=blue,
    citecolor=blue,
    urlcolor=blue,
    pdftitle={H4 Numerological Atlas: A Catalog of Standard Model Parameter Formulas},
    pdfsubject={Data paper, High Energy Physics - Phenomenology},
    pdfkeywords={Standard Model, mass formulas, golden ratio, data catalog, Coxeter group}
}

% ─── Custom commands ──────────────────────────────────────────────────────────
\newcommand{\ph}{\varphi}
\newcommand{\GeV}{\,\text{GeV}}
\newcommand{\eV}{\,\text{eV}}
\newcommand{\PDG}{\text{PDG 2024}}
\newcommand{\FLAG}{\text{FLAG 2024}}

% ─── Begin Document ───────────────────────────────────────────────────────────
\begin{document}

% ═══════════════════════════════════════════════════════════════════════════════
%  TITLE AND AUTHORS
% ═══════════════════════════════════════════════════════════════════════════════
\title{H\textsubscript{4} Numerological Atlas: A Catalog of Standard Model Parameter Formulas}
\author{[Author List TBD]\\
  \texttt{trinity-s3ai@proton.me}}
\date{May 2026}
\maketitle

% ═══════════════════════════════════════════════════════════════════════════════
%  ABSTRACT
% ═══════════════════════════════════════════════════════════════════════════════
\begin{abstract}
We present a curated catalog of 25 closed-form formulas that correlate
parameters of the Standard Model (SM) with combinations of mathematical
constants ($\varphi$, $\pi$, $e$) and invariants of the H\textsubscript{4}
Coxeter group.  The catalog is provided as a reproducible dataset
accompanied by formal verification scripts (Coq, Python), numerical
validation against Particle Data Group 2024 values, and a pre-registered
Monte-Carlo statistical test.  Mean relative agreement is $0.08\%$,
with 20 of 25 formulas achieving $<0.1\%$ error.  A pre-registered
Monte-Carlo simulation ($5\times 10^4$ trials, $2\times 10^3$ random
formulas per trial) yields $p < 10^{-4}$ against the null hypothesis
of random coincidence.  All data, code, and proofs are available on
GitHub and archived with Zenodo.
\end{abstract}

% ═══════════════════════════════════════════════════════════════════════════════
%  INTRODUCTION
% ═══════════════════════════════════════════════════════════════════════════════
\section{Introduction}

The Standard Model of particle physics is defined by approximately 25
independent parameters (masses, mixing angles, gauge couplings) whose
numerical values are determined experimentally but not predicted by the
theory itself.  Over the past five decades, numerous authors have
observed numerical coincidences between these parameters and mathematical
constants, notably the golden ratio $\varphi = (1+\sqrt{5})/2$, $\pi$,
and $e$\ \cite{Koide:1983,Connes:1996, Rivero:2006, Furey:2015}.

Such coincidences are typically reported in isolation, without systematic
statistical assessment.  The ``Trinity S$^3$AI'' program (Waves 1--19)
collects these observations into a single, machine-verified catalog and
subjects the entire set to an honest statistical test.  This paper is a
\textbf{data paper}: it documents the catalog as a reproducible dataset,
makes no claim that the formulas are ``derived'' from first principles,
and reports both the numerical agreement and its statistical significance.

\paragraph{Prior art.}  Koide's mass formula for charged leptons
achieves $9.3\,\text{ppm}$ precision\ \cite{Koide:1983} but lacks a
derivation.  Connes' noncommutative geometry program predicts the Higgs
mass from spectral action\ \cite{Connes:2006} but requires a
phenomenological $\sigma$-field correction.  The Furey--Gresnigt
Cl(8) program relates SM fermions to octonionic spinors\ \cite{Furey:2015}.
None of these programs has published a \emph{catalog} of formulas with
rigorous statistical assessment.

\paragraph{Gap.}  There is no peer-reviewed dataset that (i) collects
SM parameter formulas in one place, (ii) classifies each formula by
epistemic status (structurally motivated vs numerically fitted), (iii)
validates every entry against PDG values with documented uncertainty,
and (iv) reports a pre-registered Monte-Carlo $p$-value for the catalog
as a whole.

\paragraph{Scope.}  This paper presents exactly such a dataset.  The
accompanying negative-result paper (Wave 17)\ \cite{TrinityWave17} proves
that the H\textsubscript{4}/600-cell discrete model does \emph{not}
derive the SM.  The present paper isolates the surviving numerical
correlations and documents them as a citable data resource.


% ═══════════════════════════════════════════════════════════════════════════════
%  DATA DESCRIPTION
% ═══════════════════════════════════════════════════════════════════════════════
\section{Data Description}
\label{sec:data}

The catalog contains 25 closed-form formulas mapping H\textsubscript{4}-inspired
expressions to Standard Model observables.  Each entry includes: a unique
identifier, the mathematical expression, the target PDG 2024 value, the
computed value, relative error, $\sigma$-distance (using real experimental
uncertainty), and an epistemic classification.

\subsection{Mass scheme conventions}

The catalog uses a mixed mass scheme, following PDG and FLAG conventions:
\begin{itemize}
  \item Leptons: pole masses (on-shell at $Q^2=m^2$).
  \item Light quarks ($u$, $d$, $s$): $\overline{\text{MS}}$ at $2\GeV$.
  \item Heavy quarks ($c$, $b$): $\overline{\text{MS}}$ at self-scale.
  \item Top quark, Higgs, gauge bosons: pole masses.
  \item Neutrinos: effective Majorana (see-saw suppressed).
\end{itemize}

\subsection{Formula classification}

Every formula is tagged by epistemic status:
\begin{itemize}
  \item \textbf{(S) Structural} ($n=8$): Integer coefficients are
    H\textsubscript{4} invariants (Coxeter numbers, degrees, Lucas
    numbers).  The transcendental combination ($\varphi^a\pi^b e^c$) is
    \emph{not} derived from H\textsubscript{4} but is numerically fitted.
  \item \textbf{(NF) Numerical Fit} ($n=17$): No structural derivation;
    the entire expression is fitted to the data.
\end{itemize}
This taxonomy is adapted from the catalog audit of
Ref.~\cite{TrinityAudit}.

\subsection{The 25 core formulas}

Table~\ref{tab:formulas} presents the complete catalog.  All values are
computed with \texttt{mpmath} at 50-digit precision.  The
$\sigma$-distance is defined as
\begin{equation}
  \sigma = \frac{|\text{predicted} - \text{measured}|}{\sigma_{\text{measured}}},
\end{equation}
where $\sigma_{\text{measured}}$ is the 1-$\sigma$ experimental
uncertainty from PDG 2024, FLAG 2024, or NuFit 6.0.

% ─── Table 1: Formula catalog ─────────────────────────────────────────────────
\begin{longtable}{clccccc}
\caption{The 25 core formulas of the H\textsubscript{4} Numerological Atlas.
  Grades: SG ($\sigma<0.1$), V ($0.1\le\sigma<1$), P ($1\le\sigma<3$),
  M ($3\le\sigma<5$), F ($\sigma\ge 5$).}
\label{tab:formulas} \\
\toprule
ID & Expression & Target & Computed & Rel.Err & $\sigma$ & Grade \\
\midrule
\endfirsthead
\multicolumn{7}{c}{\tablename\ \thetable{} -- continued} \\
\toprule
ID & Expression & Target & Computed & Rel.Err & $\sigma$ & Grade \\
\midrule
\endhead
\midrule
\multicolumn{7}{r}{Continued on next page} \\
\endfoot
\bottomrule
\endlastfoot
L01 & $239e/\pi$ & 206.768 & 206.796 & 0.013\% & $6.1\times10^{3}$ & F \\
L02 & $239\varphi^{4}/\pi^{4}$ & 16.817 & 16.817 & 0.002\% & 0.29 & V \\
L03 & $549e\pi^{2}/\varphi^{3}$ & 3477.3 & 3476.99 & 0.009\% & 1.54 & P \\
Q01 & $2\varphi/7$ & 0.462 & 0.4623 & 0.064\% & 0.02 & SG \\
Q02 & $12+\varphi^{3}e^{2}$ & 43.24 & 43.30 & 0.140\% & 0.04 & SG \\
Q03 & $19\pi e^{2}/\varphi$ & 272.0 & 272.59 & 0.216\% & 0.12 & V \\
Q04 & $24\pi^{3}/e^{4}$ & 13.633 & 13.630 & 0.025\% & 0.17 & V \\
Q05 & $43+\pi/\varphi$ & 44.94 & 44.942 & 0.004\% & 0.02 & SG \\
Q05b& $127\varphi/120+30/19$ & 3.2908 & 3.2914 & 0.017\% & 0.19 & V \\
Q06 & $\pi e^{4}+6/5$ & 172.69 & 172.73 & 0.020\% & 0.12 & V \\
Q07 & $24\varphi^{2}/\pi$ & 20.00 & 20.000 & 0.002\% & 0.00 & SG \\
G01 & $36\varphi e^{2}/\pi$ & 137.036 & 137.003 & 0.024\% & $1.6\times10^{6}$ & F \\
G02 & $(\sqrt{5}-2)/2$ & 0.1179 & 0.1180 & 0.114\% & 0.13 & V \\
N01 & $8\pi/(\varphi^{5}e^{2})$ & 0.307 & 0.3067 & 0.098\% & 0.02 & SG \\
N04 & $3/\varphi^{2}\cdot180/\pi$ & 65.5$^{\circ}$ & 65.66$^{\circ}$ & 0.237\% & 0.10 & SG \\
C01 & $2\varphi^{3}e^{2}/(9\pi^{3})$ & 0.2265 & 0.2243 & 0.958\% & 4.34 & M \\
C02 & $1/(3\varphi^{2}\pi)$ & 0.04053 & 0.04053 & 0.005\% & 0.00 & SG \\
H01 & $4\varphi^{3}e^{2}$ & 125.20 & 125.20 & 0.002\% & 0.02 & SG \\
H02 & $11\varphi/20+2/3$ & 1.5577 & 1.5566 & 0.067\% & 0.75 & V \\
H03 & $4\varphi\pi/15+4/225$ & 1.3737 & 1.3733 & 0.022\% & 0.25 & V \\
$\nu$21& $(\varphi e/\pi)^{6}\times10^{-5}$ & $7.53\times10^{-5}$ & $7.53\times10^{-5}$ & 0.0003\% & 0.00 & SG \\
$\nu$31& $15\varphi^{-5}\pi^{-2}e^{-4}$ & $2.51\times10^{-3}$ & $2.51\times10^{-3}$ & 0.0004\% & 0.00 & SG \\
N21 & $\pi/(40\varphi^{2})$ & 0.0300 & 0.0300 & 0.002\% & 0.00 & SG \\
Pr  & $6\pi^{5}$ & 1836.15 & 1836.12 & 0.002\% & $3.1\times10^{5}$ & F \\
S$\nu$& $8\varphi^{-6}\pi^{-5}e^{6}\times0.1$ & 0.0588 & 0.05877 & 0.045\% & 0.00 & SG \\
\end{longtable}

\subsection{The ultra-precision trap}

Three formulas receive ``Fail'' grades on $\sigma$-distance not because
their relative errors are large (all three have $<0.03\%$ error) but
because the corresponding experimental measurements are extraordinarily
precise:
\begin{itemize}
  \item $1/\alpha(0) = 137.035\,999\,084(21)$ (relative uncertainty
    $1.5\times10^{-10}$)
  \item $m_{p}/m_{e} = 1836.152\,673\,43(11)$ (relative uncertainty
    $6.0\times10^{-11}$)
  \item $m_{\mu}/m_{e} = 206.768\,283\,0(46)$ (relative uncertainty
    $2.2\times10^{-9}$)
\end{itemize}
Any catalog claiming to ``derive'' these constants must match
experimental precision ($\sim 10^{-10}$), not merely approximate them
($\sim 10^{-4}$).  The honest conclusion is that the catalog
\emph{correlates} these constants but does not \emph{derive} them.

% ═══════════════════════════════════════════════════════════════════════════════
%  METHODS
% ═══════════════════════════════════════════════════════════════════════════════
\section{Methods}
\label{sec:methods}

\subsection{Formula discovery}

The formulas were discovered by searching a space of expressions of the
form
\begin{equation}
  f = C \, \varphi^{a} \, \pi^{b} \, e^{c},
\end{equation}
and four related templates (additive, ratio, two-term), where $C$ is a
rational prefactor $p/q$ with $p\le 1000$, $q\le 100$, or an explicit
H\textsubscript{4} invariant $\{1,2,7,11,12,19,20,29,30,120,240,480\}$,
and exponents $a,b,c\in\{-6,\dots,6\}$.  The total search space is
$\approx 6.7\times10^{8}$ distinct expressions per template.

The discovery process was \emph{not} automated: candidate formulas were
proposed based on H\textsubscript{4} invariants and tested numerically.
This manual search is a source of potential bias (see Limitations).

\subsection{Numerical validation}

Every formula is evaluated with \texttt{mpmath}\ \cite{mpmath} at
50-digit precision and compared to PDG 2024, FLAG 2024, or NuFit 6.0
targets.  The validation script \texttt{validate\_v4.py} is part of the
dataset.

\subsection{Statistical test}

We performed a pre-registered Monte-Carlo simulation to test the null
hypothesis that the catalog is indistinguishable from random search:
\begin{enumerate}
  \item \textbf{Pre-registration} (2026-05-22): search space, five
    formula templates, and PDG targets were fixed before running.
  \item \textbf{Simulation}: $5\times10^{4}$ trials, each drawing
    $2\times10^{3}$ random formulas from the search space.
  \item \textbf{Metric}: for each trial, the best relative error per
    observable was recorded.
  \item \textbf{Comparison}: Trinity's error vector was compared to the
    distribution of random best errors.
\end{enumerate}


% ═══════════════════════════════════════════════════════════════════════════════
%  DATA RECORDS
% ═══════════════════════════════════════════════════════════════════════════════
\section{Data Records}
\label{sec:records}

The dataset is organized as follows:

\begin{itemize}
  \item \texttt{FORMULAS.md} — Human-readable catalog with full
    derivations, classification justifications, and cross-references.
  \item \texttt{validate\_v4.py} — Python validation script (mpmath,
    50-digit precision).
  \item \texttt{scripts/honest\_phenomenology\_v18.py} — Monte-Carlo
    p-value engine (vectorized, numpy).
  \item \texttt{derivations/catalog\_audit/} — Audit reports:
    \begin{itemize}
      \item \texttt{sigma\_ranking.md} — $\sigma$-distance ranking
      \item \texttt{honest\_pvalue\_report\_v18.md} — Full MC report
      \item \texttt{wave18\_mc\_results.json} — Raw MC results
    \end{itemize}
  \item \texttt{proofs/trinity/} — Coq formal proofs (1,772 \texttt{Qed},
    0 \texttt{Admitted}, 77 Axioms).
  \item \texttt{derivations/lean\_port/} — Lean 4 port (9 modules).
\end{itemize}

\subsection{JSON schema}

Each formula entry follows this schema:
\begin{verbatim}
{
  "id": "Q07",
  "expression": "24*PHI**2/PI",
  "physical_meaning": "m_s/m_d (2 GeV MS-bar)",
  "target_value": 20.00,
  "target_uncertainty": 0.70,
  "computed_value": 20.0003,
  "rel_error_percent": 0.0015,
  "sigma_distance": 0.00,
  "classification": "S",
  "source_file": "proofs/trinity/Catalog42.v"
}
\end{verbatim}

% ═══════════════════════════════════════════════════════════════════════════════
%  TECHNICAL VALIDATION
% ═══════════════════════════════════════════════════════════════════════════════
\section{Technical Validation}
\label{sec:validation}

\subsection{Monte-Carlo results}

Table~\ref{tab:mc} summarizes the Monte-Carlo comparison.  The Trinity
catalog achieves significantly smaller mean relative errors than random
sampling ($p<10^{-4}$).  The number of high-precision hits ($<0.1\%$
error) is also highly significant ($p<10^{-4}$).

\begin{table}[h]
\centering
\caption{Monte-Carlo comparison: Trinity vs random search
  ($5\times10^{4}$ trials, $2\times10^{3}$ formulas/trial).}
\label{tab:mc}
\begin{tabular}{lcccc}
\toprule
Metric & Trinity & Random (median) & Random (mean$\pm$std) & $p$-value \\
\midrule
Mean rel.\ error (\%) & 0.0835 & 0.594 & $0.624\pm0.185$ & $<10^{-4}$ \\
Hits $<0.1\%$ & 20/25 & 5/25 & $5.2\pm2.0$ & $<10^{-4}$ \\
Hits $<1.0\%$ & 25/25 & 21/25 & $21.1\pm1.6$ & 0.0061 \\
\bottomrule
\end{tabular}
\end{table}

\subsection{Per-formula significance}

Only 10 of the 25 formulas are individually significant at $p<0.05$;
the remaining 15 are consistent with random search when tested in
isolation.  The catalog's strength is therefore \emph{collective}, not
individual.  The significant formulas are: L02, L03, Q05, Q07, C02, H01,
$\nu$21, $\nu$31, N21, and Pr.

\subsection{Limitations}

\begin{enumerate}
  \item \textbf{Search space uncertainty}: The true space searched during
    development may differ from the pre-registered $\sim$$10^{8}$ space.
  \item \textbf{Template limitations}: Only five templates were tested;
    the original search may have used additional forms.
  \item \textbf{Manual discovery bias}: Formulas were proposed manually,
    not by exhaustive search, introducing potential confirmation bias.
  \item \textbf{No Bonferroni correction}: Aggregate $p$-values are
    uncorrected for multiple testing.
  \item \textbf{Upper-limit observables}: $\sum m_{\nu}$ is a 95\% CL
    upper limit, not a Gaussian measurement; its $\sigma$-distance is
    less meaningful.
\end{enumerate}

% ═══════════════════════════════════════════════════════════════════════════════
%  USAGE NOTES
% ═══════════════════════════════════════════════════════════════════════════════
\section{Usage Notes}

\paragraph{Reproduce the validation.}
\begin{verbatim}
  $ python3 validate_v4.py
\end{verbatim}
Expected output: 25 formulas validated, mean relative error $0.08\%$,
100\% pass rate at $<1\%$ threshold.

\paragraph{Reproduce the Monte-Carlo test.}
\begin{verbatim}
  $ python3 scripts/honest_phenomenology_v18.py
\end{verbatim}
Expected runtime: $\sim$2 minutes (50,000 trials, 2,000 formulas/trial).

\paragraph{Extend the catalog.}
Add new formulas to \texttt{validate\_v4.py} and \texttt{FORMULAS.md},
update PDG targets, and re-run validation.  The MC engine accepts
arbitrary formula templates via the \texttt{TEMPLATE\_REGISTRY}.

% ═══════════════════════════════════════════════════════════════════════════════
%  DATA AVAILABILITY
% ═══════════════════════════════════════════════════════════════════════════════
\section{Data Availability}

All data, code, and formal proofs are available at:
\begin{itemize}
  \item \textbf{GitHub}: \url{https://github.com/gHashTag/trinity-s3ai}
  \item \textbf{Zenodo}: DOI 10.5281/zenodo.XXXXXX (pending publication)
  \item \textbf{License}: MIT (code) / CC-BY-4.0 (data)
\end{itemize}

The repository includes:
\begin{itemize}
  \item 79 Coq proof files (1,772 \texttt{Qed}, 0 \texttt{Admitted})
  \item 9 Lean 4 modules
  \item Python validation and MC scripts
  \item LaTeX source of this paper
  \item JSON data files
\end{itemize}

% ═══════════════════════════════════════════════════════════════════════════════
%  ACKNOWLEDGEMENTS
% ═══════════════════════════════════════════════════════════════════════════════
\section*{Acknowledgements}

We thank the Coq, Lean, and mpmath communities for their open-source
tools.  This work was supported entirely by curiosity.

% ═══════════════════════════════════════════════════════════════════════════════
%  REFERENCES
% ═══════════════════════════════════════════════════════════════════════════════
\begin{thebibliography}{99}

\bibitem{Koide:1983}
Y. Koide, ``New View of Quark and Lepton Mass Matrix,''\n  \textit{Phys.\ Rev.\ D} \textbf{28} (1983) 252.

\bibitem{Connes:1996}
A. Connes, ``Gravity Coupled with Matter and Foundation of
  Noncommutative Geometry,'' \textit{Commun.\ Math.\ Phys.}\ \textbf{182}
  (1996) 155, arXiv:hep-th/9603053.

\bibitem{Connes:2006}
A. Connes and A. H. Chamseddine, ``Inner Fluctuations of the Spectral
  Action,'' \textit{J.\ Geom.\ Phys.}\ \textbf{57} (2006) 1.

\bibitem{Rivero:2006}
A. Rivero, ``A Lepton Mass Formula,'' arXiv:hep-ph/0605220 (2006).

\bibitem{Furey:2015}
C. Furey, ``Standard Model Physics from an Algebra?''
  arXiv:1611.09182 [hep-th] (2015).

\bibitem{Gresnigt:2021}
N. Gresnigt, ``On the Octonionic Standard Model,'' arXiv:2107.00569
  [hep-th] (2021).

\bibitem{PDG:2024}
Particle Data Group, ``Review of Particle Physics,''
  \textit{Phys.\ Rev.\ D} \textbf{110} (2024) 030001.

\bibitem{FLAG:2024}
Flavour Lattice Averaging Group (FLAG),
  ``Review of Lattice Results concerning Low-Energy Particle Physics,''
  arXiv:2402.03554 [hep-lat] (2024).

\bibitem{NuFit:2024}
NuFit Collaboration, ``NuFit 6.0,'' \url{www.nu-fit.org} (2024).

\bibitem{mpmath}
F. Johansson et al., ``mpmath: a Python library for arbitrary-precision
  floating-point arithmetic,'' \url{mpmath.org} (version 1.3.0).

\bibitem{TrinityWave17}
Trinity S$^3$AI Collaboration, ``An Active Boundary-Mapping Research Program in
  Geometric Unification,'' arXiv:XXXX.XXXXX [hep-th] (2026).

\bibitem{TrinityAudit}
Trinity S$^3$AI Collaboration, ``Catalog Audit Report,''
  \url{github.com/gHashTag/trinity-s3ai/blob/main/derivations/catalog_audit/audit_report.md}.

\bibitem{Brannen:2012}
C. Brannen, ``The Lepton Masses,'' arXiv:1202.4079 [physics.gen-ph] (2012).

\bibitem{Sumino:2009}
Y. Sumino, ``Family Gauge Bosons and Koide's Mass Formula,''
  arXiv:0909.4429 [hep-ph] (2009).

\bibitem{Dechant:2013}
P.-P. Dechant, ``Platonic Solids Generate their Four-Dimensional
  Analogues,'' \textit{Acta Cryst.}\ A\textbf{69} (2013) 592.

\bibitem{Lisi:2007}
G. Lisi, ``An Exceptionally Simple Theory of Everything,''
  arXiv:0711.0770 [hep-th] (2007).

\bibitem{Distler:2009}
J. Distler and S. Garibaldi, ``There is no `Theory of Everything' inside
  E$_8$,'' \textit{Commun.\ Math.\ Phys.}\ \textbf{298} (2010) 419.

\bibitem{Ramond:2001}
P. Ramond, ``Exceptional Groups and Physics,''
  arXiv:hep-th/0110032 (2001).

\bibitem{Gursey:1975}
F. G\"{u}rsey and P. Sikivie, ``E$_6$ as a Universal Gauge Group,''
  \textit{Phys.\ Rev.\ Lett.}\ \textbf{36} (1976) 775.

\bibitem{PLOS:2024}
PLOS ONE, ``Data Availability and Transparency,''\n  \url{journals.plos.org/plosone/s/data-availability} (2024).

\end{thebibliography}

\end{document}
